\documentclass[12pt]{article}
\usepackage[utf8]{inputenc}
\usepackage[margin=0.75in]{geometry}

\title{CSC110 Fall 2026 Assignment 1: Written Questions}
\author{Idzudin Saffuan}
\date{\today}

\begin{document}
\maketitle

\section*{Part 1: Interpreting test results}

% TIP: In LaTeX, the underscore (_) is a special character, so if you want to use it
% in normal text, you have to put a backslash in front of it. E.g., a1\_part2.py,
% not a1_part2.py.

\begin{enumerate}
\item[1.]
test\_calculation\_one\_group\_one\_colour PASSED

test\_calculation\_many\_groups\_same\_colours FAILED

test\_calculation\_many\_groups\_different\_colours FAILED

\item[2.]
Error 1 : colours[0][2] instead of colour[2]

- The error causes the function to read the g value for the first colour in the list.

- Rather, the g value should be read from within each of the colour given to the function.

Error 2 : return total / len(colours) instead of return total

- The error causes the function to return the average of the brightness among the colours.

- Rather, the function should return the total brightness of the colours.

\item[3.]
For a list with one colour, Error 1 still returns the correct value as reading from colours[0] refers to the same colour.

Error 2 still returns the correct value as finding the average of brightness across a single colour is equal to the brightness of the colour itself.

\end{enumerate}

\section*{Part 4: Reasoning about Arguments}

\begin{enumerate}
\item[Claim 1.]
Claim 1 and Paul's explanation is correct given that the original value of B is not 255.

\item[Claim 2.]
Both Claim 2 and Paul's reasoning is incorrect as the possible colour balance is from -255 to 255.

In a collection of 512 valid RGB colours, it is not necessarry that two colours have the same colour balance.

\item[Claim 3.]
Claim 3 is correct but Paul's explanation is inaccurate.

For cases where the colour balance is negative, the property that
$ac > bc$ is incorrect for $a > b$.

The correct explanation is that $a^2 > b^2$ for $a > b$ hence the greater the colour balance, the greater the colour imbalance.
\end{enumerate}

\end{document}
